% !TeX root = ../prompt-engineering-guide-for-beginners.tex

%%%%%%%%%%%%%%%%%%%%%%%%%%%%%%%%%%
%% 第三章：基本概念与模型设置
%%%%%%%%%%%%%%%%%%%%%%%%%%%%%%%%%%
\chapter{基本概念：理解大语言模型的工作方式}
\label{ch:basic-concepts}

在动手书写更复杂的提示词之前，我们先花几分钟了解一些大语言模型（LLM）的基本工作概念。
放心，这些概念并不艰涩，但掌握它们能帮助你理解 AI 底层的行为逻辑，从而写出更精炼、更有效的提示词。

\section{大语言模型的基本原理}

上一章我们提到，大语言模型本质上是一个「超级文字接龙选手」。这一节我们在此基础上稍微深入一点，介绍三个核心技术概念。

\subsection*{Token（令牌）}
\label{sec:token}

AI 不是像人类一样逐字逐句地阅读你输入的内容，而是先将文本切分成一个个更小的语义单元，这些单元在技术上叫做 Token。

对于英文来说，一个 Token 大约相当于一个完整的单词或单词的一部分（例如 "hamburger" 在某些模型引擎中会被 BPE（字节对编码）算法切分成 "ham" 和 "burger"）；对于中文来说，由于信息密度大，一个汉字通常对应 1 到 2 个 Token，具体取决于不同厂商底层模型训练时使用的词表（Vocabulary）切分与编码规则。

举个具体的例子来感受一下：「今天上海天气蛮好的」这句话大约会被切分成 7 到 12 个 Token（取决于使用的具体模型），而对应的英文 "The weather in Shanghai is great today" 大约需要 8 个 Token。各家模型厂商通常都提供了在线 Tokenizer 工具，你可以亲自输入一段文字，直观体验 Token 的切分过程。

\begin{tipbox}
为什么要深度了解 Token？因为在模型眼底，世界上没有“字”的概念，只有一长串对应的 Token ID 数字。模型对代码或长串数字的拆分往往违反人类直觉（例如把“3.1415”切成“3.14”和“15”），这就是产生弱智计算错误的根本原因。同时，所有大模型 API 接口计费、输入输出额度限制，全部都是以 Token 数量为标准计算的。
\end{tipbox}

\subsection*{上下文窗口（Context Window）}
\label{sec:context-window}

上下文窗口是 AI 在一次对话中能「看到」的信息总量上限，同样是用 Token 数量来衡量的。

你可以把它想象成为一张物理的「工作台」：工作台的面积是固定的，上面能摆放的文件也有限。只要不超过这个范围，AI 都能综合这些信息给出回答；但一旦你提供的信息超出范围，AI 就只能把相对「陈旧」的信息挤出工作台，结果就是它「忘记」了你之前的交代。

早期的模型上下文可能只有几千个 Token（约等于几页纸的内容），而目前的先进模型（例如 Kimi、Claude 3 等）有些已经支持几十万甚至上百万 Token（足以一次性塞入整整一本书的内容）。

\subsection*{补全（Completion）}
\label{sec:completion}

从技术底层来看，AI 进行信息生成的过程叫做「补全」。你在对话框内提交一段文字（即提示词），AI 所做的主要工作是根据统计概率，紧接着这段文字向下「续写」生成它认为最合理的后续内容。

即使你看起来是在以一问一答的方式发出问题，AI 其实也不在使用传统的思维去「思考」和「回答」，而是在计算这段对话接下来最可能出现怎样的后文。这就是为什么提示词起到了四两拨千斤的作用：你给的开头和背景哪怕只有微妙的区别，它续写演化的方向都会截然不同。

\section{大语言模型常规设置（LLM Settings）}

许多面向开发者和深度用户的 AI 工具（例如各种调用 API 的客户端、Playground 等）都会开放一些参数设置面板，允许你自定义或者干预模型的输出行为。掌握这些参数，你就能让 AI 变得按照你的特殊想法进行输出。

\subsection*{Temperature（温度）}
\label{sec:temperature}

Temperature 是出镜率最高的核心控制参数，它主导着大模型输出的「确定性与随机性」。数值一般从 0 开始。

\begin{itemize}
  \item \textbf{低温度值}（如 0.1 到 0.3）：剥夺 AI 的大部分发散能力。它的回答将变得非常确定、聚焦且极其稳定。该模式适合事实精准查询、结构化数据提取、代码生成、逻辑严密的数学计算等「只求准不求怪」的任务。
  \item \textbf{高温度值}（如 0.7 到 1.0 乃至更高）：解锁 AI 的天马行空。它的回答会充满创意与意想不到的表达方式，但相伴的风险就是更容易出现胡编乱造的「幻觉」。适合用在头脑风暴、虚构文学创作或是灵感启发的场景。
\end{itemize}

来感受一下 Temperature 对同一个问题的影响。假如你问 AI「用一句话形容上海」：

\begin{itemize}
  \item \textbf{低温度（0.2）}可能回答：「上海是中国最大的经济中心城市。」
  \item \textbf{高温度（0.9）}可能回答：「上海是一座让人在弄堂里吃着葱油拌面、却仰头就能看到陆家嘴天际线的魔法城市。」
\end{itemize}

同一个 AI，同一个问题，仅仅因为温度参数不同，输出的「画风」就判若两人。

\begin{tipbox}
口诀：追求靠谱选低温，想要创意推高温。
\end{tipbox}

\subsection*{Top-p（核采样）}

Top-p 是另一种常见的、用于调控输出概率多样性的参数。虽然它的技术原理与 Temperature 不同，但表现在结果上的作用极其相似。

Top-p 的设定值位于 0 到 1 之间。值越小，模型在选下一个词时就越拘谨（它只会从少数几个概率最高、最常出现的词汇中去进行选择）；值越大，那些原本冷门或平庸的词，也会有被选中的机会。

\begin{tipbox}
通常情况下，开发者们约定俗成的一条原则是：如果你想调整结果的丰富度，\textbf{只调整 Temperature 或只调整 Top-p 中的一个即可}，这两者没必要同时改动，否则容易让产出失控。
\end{tipbox}

\subsection*{Max Tokens（最大令牌数）}

此参数是一道红线，强制限定了 AI 单次回答时能生成的最大长度（Token 总量）。如果用户只想要一句简短的口号，把这个值压低，就能避免部分「过度热情」的大模型喋喋不休。

假设你将 Max Tokens 设为 100，当生成内容触及第 100 个 Token 时，不管句子是否讲完、标点符号有没有闭合，模型都会像被断掉电源一样立刻停止。这就要求你在实际应用中，不能一味图短而设得过于极致，以免最终得到被生硬截断的半成品回答。

\subsection*{Stop Sequences（停止序列）}

用户可以向系统中传入一个或多个特定的字符串列表（比如某种标记符号或某些句子），模型一旦在它的接龙过程中生成了与这个列表相匹配的内容，生成进程就会立即停下。

以多角色剧本生成为例，假如你希望每次生成一人份的台词，你可以将「李雷：」设定为停止序列。这样，系统在完成一段输出并准备开始「李雷：」时，就会被强行中断。如此一来，你将拿到单独一轮回合的干净文本，而不至于长篇大论。这个功能在代码提取或者 API 设计中常被作为截断垃圾文本的操作，虽然我们在日常文字聊天时鲜少触及。

\subsection*{Frequency / Presence Penalty（频率惩罚与存在惩罚）}

很多时候你可能遇到 AI 开始循环重复特定词语或陷入「车轱辘话」中，这两个参数就是应对该烦恼的克星。

Frequency Penalty（频率惩罚）的核心规则是按出现次数递减：模型在单次生成中每用过一次某词汇，该词再次出场的概率就在递减，重复越多打压越狠。
而 Presence Penalty（存在惩罚）奉行的则是「一视同仁的打压机制」：无论某个词汇出现了 1 次还是 10 次，只要「存在于本文中」，再次调用的门槛就会被施加一次固定的难度。

如果在实践中发现模型像口头禅一样频繁抛出「总而言之」、「另外一方面」（车轱辘话），你可以适当地调高这两个惩罚值，驱使模型调用更多新的词汇；代价则是如果调得过高，内容就会显得生硬和词不达意，因为它可能会为了避开常用词语而被「逼着」抛出生僻字词。

\subsection*{不同任务怎么调？}

下面这张速查表可以帮你快速决定参数设置：

\begin{center}
\begin{tabular}{lcc}
  \toprule
  \textbf{任务类型} & \textbf{Temperature} & \textbf{建议} \\
  \midrule
  事实问答、数据分析 & 0.1\textasciitilde 0.3 & 追求准确 \\
  邮件、报告、翻译    & 0.3\textasciitilde 0.5 & 稳定可靠 \\
  文案创作、头脑风暴  & 0.7\textasciitilde 0.9 & 鼓励创意 \\
  诗歌、小说、自由创作 & 0.9\textasciitilde 1.0 & 最大随机性 \\
  \bottomrule
\end{tabular}
\end{center}

如果你实在不知道怎么设置，保持默认值就好。大多数 AI 工具的默认参数已经针对通用场景做了优化。

\section{系统提示词（System Prompt）}
\label{sec:system-prompt}

许多现代大语言模型的应用（例如 ChatGPT、Claude 以及其他大语言模型生态软件如 Cherry Studio、AnythingLLM）都会专门提供一个功能：「系统提示词 / 系统预设」设定。

系统提示词的功能，其实是一道具有更高权限的背景环境预设墙。在它身上可以完成以下指令动作：

\begin{itemize}
  \item 奠定专业基础与人设：例如「你是一位有着 20 年经验的心胸外科医生，接下来的对话中要保持高度的同理心。」
  \item 控制核心风格：例如「始终只用严谨得体的规范中文输出，禁止使用网络缩写词及任何外语单词的混入使用。」
  \item 设置铁律与工作边界：例如「但凡遇到任何越界问题，你都必须立刻回复『很抱歉我不具备相关权限』，严禁任何擅自的脑补。」
\end{itemize}

你提交的这些系统预设，并不会夹杂在你平时的对话问答里。它隐匿于前端幕布之后，却是影响和重塑整个对话过程模型最终走向的最高统帅力量。

\begin{tipbox}
即使你使用的是没有专门设置入口的直接对话框，通过在第一次聊天发送的长提示词中直接带上一句「在接下来的整个会话窗口中，请将自己定位为一位富有经验的高级财务顾问」，这也是通过模拟系统提示词打下的工作基调。
\end{tipbox}

\section{多轮对话与上下文管理}
\label{sec:multi-turn}

绝大多数场景中，大家感受不到 AI 在和自己互动沟通中产生断层的违和感，因为它能承接上文的语境。这是由于 AI 在幕后做了一件体力活，它把从对话开启至今你发出的每一条记录和它做出的每一句回答全部放在了自己的上下文窗口中。

遗憾的是，这种记忆方式绝非无限容量。正如前面小节解释的，我们已经知道了它是把记忆放置在「一定尺寸的物理工作台」上的。如果你们在这间聊天室中不断对答，产生的内容在字数上限度超过了当前平台所规定的窗口极限尺寸时，系统便会自动启动机制，将最先产生的对话（通常指的是第一句话和后续的早期互动内容）强行抹去，来为新生内容让路。这就造成了我们在连续对话数日后发现，AI “失去记忆”，完全忘记了当初角色扮演设定。

面对这一大模型的底层短板，目前有以下几点有效的提示工程应对小窍门：

\begin{itemize}
  \item \textbf{尽早开启新对话分支（Reset Session）}：如果你目前的交流与先前的分析毫无关系，请务必「另起炉灶」清空当前记录（重启 Chat Session）。这样不但可以还 AI 一块空白的“无污染算力画板”，更能避免由于上下文过于庞大导致 AI 注意力分散和 API 额度浪费。
  \item \textbf{利用阶段性回溯进行状态压缩（State Compression）}：如果是跨越长文本的写作或分析过程。在感到上下文快吃满前，请让 AI 执行：「基于以上沟通进展，请输出一份结构化总结」。随后你可以开启新对话，并把这份“高纯度压缩包”作为新环境的起始状态（Prompt），直接跨过冗长无用的聊天记录。
  \item \textbf{重申关键铁律信息}：面对开始偏离轨道的 AI，与其在对话框内去用自然语言与它反复争辩、指责它忘记旧规则（这反而会在上下文中注入更多污染性的 Token），不如直接在新的命令前再叠加一次最高权重的提示内容：「请再次使用开局定义的规则去作回复」。
  \item \textbf{对抗“迷失在中间（Lost in the Middle）”效应}：多项论文研究表明，当上下文达到数万 Token 时，模型抓取顶端和尾部信息的注意力分数始终保持最高，而中间段落的信息极易被忽略（即所谓的 U 型注意力曲线）。因此如果有影响存亡的规矩、或者极为核心的输出格式参考，请务必将其排布在整个 Prompt 的最顶部（作为 System 预设）或者最末尾的总结请求处。
\end{itemize}

系统参数、基础工作概念和记忆法则这三块基础知识，是你后续运用各种提示技巧时最重要的底座支撑。

\section{补充：推理模型与「深度思考」}
\label{sec:reasoning-models}

近一两年，你可能听说过一类特别的模型，如 OpenAI 的 o1/o3 系列、DeepSeek R1 等。它们被称为「推理模型」或「深度思考模型」，和普通的对话模型有一个关键区别。

普通对话模型在收到你的提示词后，会直接开始逐个生成回答的每一个 Token。而推理模型会在生成最终回答之前，先在内部进行一段较长的「思考过程」（Think / Chain of Thought），把问题拆解、验证、反复推敲之后，再给出一个更加深思熟虑的答案。

用一个直观的类比来理解：普通对话模型像是一个脱口而出的即兴演讲者，话到嘴边就说出来；而推理模型更像是一位考场上的考生，先在草稿纸上列出思路、验算一遍，确认无误后才把答案誊写到答题卡上。

\textbf{来看一个具体的例子。}假设你问 AI 这样一道逻辑题：

\begin{promptbox}
一个房间里有 3 个开关，分别控制隔壁房间的 3 盏灯。你只能进入隔壁房间一次。请问如何确定每个开关分别对应哪盏灯？
\end{promptbox}

如果使用普通对话模型，它可能会给出一个笼统的回答，甚至直接跳到结论而漏掉关键的推理步骤。但推理模型会在内部展开一段完整的思考链（你通常可以在界面上看到折叠的「思考过程」）：先分析约束条件（只能进入一次），再逐步推导出利用灯泡温度作为额外信息的策略，最终给出严密的三步解法。

\textbf{再看一个数学场景。}你让 AI 计算：

\begin{promptbox}
一个水池有两个进水管和一个出水管。A 管单独注满需要 6 小时，B 管单独注满需要 8 小时，出水管单独放完需要 12 小时。三管同时打开，多久能注满水池？
\end{promptbox}

普通模型可能直接给出一个数字却在中间步骤中犯运算错误。推理模型则会将问题拆解为分数运算，列出 $\frac{1}{6} + \frac{1}{8} - \frac{1}{12}$ 的计算过程，通分、化简，最终得出正确答案，并且每一步都可以被验证。

\textbf{什么时候用推理模型？}

\begin{itemize}
  \item 数学计算、逻辑推理等需要严密推导的任务
  \item 复杂的代码编写和调试（例如涉及多个模块的架构设计或难以定位的 Bug 排查）
  \item 需要多步骤分析和决策的问题（例如商业方案的利弊权衡、多约束条件下的方案选择）
  \item 需要对自身回答进行自我检验和纠错的场景
\end{itemize}

\textbf{什么时候用普通模型就够了？}

\begin{itemize}
  \item 日常对话、文案写作、翻译
  \item 简单的问答和信息整理
  \item 对响应速度要求较高的场景（推理模型由于需要经历「思考」阶段，通常响应更慢，消耗的 Token 也更多，费用相应更高）
\end{itemize}

\begin{tipbox}
使用推理模型时，提示词反而不需要太复杂。你不必像普通模型那样加上「请一步一步思考」之类的引导语，因为它本身就会自动进行深度推理。只需要把问题描述清楚、把约束条件列全即可。过多的格式要求或引导性话术反而可能干扰模型自身的推理节奏。
\end{tipbox}

理解了这些基础概念之后，接下来我们将开始学习如何组装提示词的各个要素。
